\section{Thermodynamic limits and the boundary-flux obstruction}
\label{therm4-section}

Throughout this section $a>0$, $g^2>0$, $m\ge0$, and $r\ge0$ are
fixed.  Energies use the raw Hamiltonian \eqref{lat-hamiltonian}.
We prove a vacuum thermodynamic limit, explicit fixed-cell relative
bounds, and an open-boundary thermodynamic limit for periodic sources.
For periodic source boundary conditions we give an exact reduction to
boundary-flux sectors and identify the additional comparison still
needed.  All gauge rotors remain untruncated.

\subsection{Open graphs and the cost of cutting a link}

For a finite rectangular set $Y\subset\mathbb Z^3$, retain its internal
nearest-neighbour links and internal plaquettes.  The open Hamiltonian
$H_Y^{\rm op}$ has these gauge terms and the fermion matrix
\begin{align}
 h_Y(U)&=m\beta+\sum_{e=(x,x+e_j)\subset Y}h_e(U),\nonumber\\
 h_e(U)&=\frac{\alpha_j}{2ia}
 \bigl(\overline U_e|x\rangle\langle y|-U_e|y\rangle\langle x|\bigr)
 +\frac{r\beta}{2a}
 (|x\rangle-U_e|y\rangle)(\langle x|-\overline U_e\langle y|).
 \label{therm4-open-matrix}
\end{align}
The one-particle spatial matrices act trivially on spin.  In particular,
the Wilson diagonal at $x$ is $m+r\deg_Y(x)/(2a)$.
It is not the periodic diagonal $m+3r/a$ at a boundary.
The Wilson mass remains at least $m$ because every link contribution
to it is positive.  The bound $\|h_Y(U)\|\le h_*$ from
\eqref{lat-mass-positive} remains valid when the degree is at most six.

Write $e_{\rm op}(Y)$ for the unrestricted auxiliary ground energy.
For an integer source $\eta:Y\to\mathbb Z$, write
$E_Y^{\rm op}(\eta)$ for the ground energy subject to
\[
 \operatorname{div}_Y E+q=\eta.
\]
This is the zero-exterior-flux convention; there are no exterior link
rotors in $H_Y^{\rm op}$.  The sector is nonempty precisely when
$|\sum_Y\eta|\le2|Y|$, for connected $Y$.
The determinant and gauge-bridge proof of Theorem~\ref{lat-main} applies
to this graph without alteration: the incidence matrix replaces the
periodic divergence, $M_Y(U)\ge m$, and the internal-link heat kernels
are positive.  Its global-angle argument likewise gives
\begin{equation}
 e_{\rm op}(Y)=E_Y^{\rm op}(0),\qquad
 E_Y^{\rm op}(\eta)\ge e_{\rm op}(Y)+m\left|\sum_Y\eta\right|.
 \label{therm4-open-comparison}
\end{equation}
Self-adjointness, compact resolvent, and finite heat traces follow as
before, including a graph with no links.  Uniform extensive bounds are
\begin{equation}
 -2h_*|Y|\le e_{\rm op}(Y)\le\frac{3}{g^2a}|Y|.
 \label{therm4-open-extensive}
\end{equation}
For the upper bound take all internal electric fields zero and two
negative-$\beta$ spinors per site.  The negative fermionic expectation
may be discarded and there are at most $3|Y|$ plaquettes.

\begin{lemma}[Bounded cutting cost with explicit rotor factors]
\label{therm4-cut}
Let a finite cubic lattice graph be partitioned into rectangular
open blocks, and let $\mathcal C$ contain all $n_c$ original links
not internal to those open blocks. This includes periodic closing
links, even when both endpoints belong to the same selected block.
On the original auxiliary Hilbert space, define
\[
 H_{\rm cut}=\sum_B H_B^{\rm op}
                 +\frac{g^2}{2a}\sum_{e\in\mathcal C}E_e^2.
\]
The last summands act on independent copies of $L^2(U(1))$.
Then
\begin{equation}
 H-H_{\rm cut}=V_{\mathcal C},\qquad
 \|V_{\mathcal C}\|\le Bn_c,\qquad
 B=\frac{4r+2+8/g^2}{a}.
 \label{therm4-cut-norm}
\end{equation}
Consequently the unrestricted ground energies differ by at most
$Bn_c$, and $\inf\sigma(H_{\rm cut})=\sum_B e_{\rm op}(B)$.
The same comparison applies when adding periodic closing links to an
open cube, after adjoining their free rotor factors.
\end{lemma}
\begin{proof}
The eight-dimensional endpoint matrix in \eqref{therm4-open-matrix}
has norm at most $r/a+1/(2a)$.  It anticommutes with $C_5$.
Its opposite eigenvalue pairs imply
\[
 \|d\Gamma(h_e)\|=\tfrac12\operatorname{Tr}|h_e|
 \le4\|h_e\|\le(4r+2)/a.
\]
A plaquette multiplier has norm at most $2/(g^2a)$, and a link meets
at most four plaquettes.  Remove each cut link's fermionic term and
every plaquette touching a cut link.  Counting a removed plaquette
more than once only enlarges the stated norm bound.
The unbounded electric term on a cut link is retained, not estimated
in norm or discarded.  Its independent rotor has ground energy zero.
Even fermion operators on distinct blocks commute under the canonical
graded Fock-space tensor decomposition.  Thus the block ground
energies and the free-rotor minima add.  The bounded-perturbation
variational principle proves the claim.
\end{proof}

\subsection{The physical vacuum thermodynamic limit}

Let $Q_n=\{0,\ldots,n-1\}^3$, $e_{\rm op}(n)=e_{\rm op}(Q_n)$, and
let $E_0^{\rm per}(n)$ be the physical vacuum energy on the $n^3$ torus.

\begin{theorem}[Vacuum limit with a finite-size estimate]
\label{therm4-vacuum-limit}
The finite limit
\begin{equation}
 e_{\rm vac}=\lim_{n\to\infty}\frac{e_{\rm op}(n)}{n^3}
            =\lim_{n\to\infty}\frac{E_0^{\rm per}(n)}{n^3}
 \label{therm4-vacuum-limit-eq}
\end{equation}
exists.  For every integer $\ell\ge3$,
\begin{equation}
 \left|e_{\rm vac}-\frac{e_{\rm op}(\ell)}{\ell^3}\right|
 \le\frac{3B}{\ell},\qquad
 \left|e_{\rm vac}-\frac{E_0^{\rm per}(\ell)}{\ell^3}\right|
 \le\frac{6B}{\ell}.
 \label{therm4-vacuum-rate}
\end{equation}
Physical energy per spatial volume is $e_{\rm vac}/a^3$.
\end{theorem}
\begin{proof}
For $n=k\ell$, cutting the open cube into $k^3$ cubes removes
$3(k-1)n^2\le3n^3/\ell$ links.  Lemma~\ref{therm4-cut} gives
\[
 |e_{\rm op}(n)-k^3e_{\rm op}(\ell)|\le3Bn^3/\ell.
\]
For arbitrary $n$, subdivide each coordinate into intervals of length
$\ell$ and one possible remainder.  The remaining rectangular boxes
contain at most $3\ell n^2$ sites.  Their energies are bounded in
absolute value by a fixed constant times that site count, by
\eqref{therm4-open-extensive}.  The number of cut links is at most
$3n^3/\ell$.  Therefore, at every fixed $\ell$,
\[
 \frac{e_{\rm op}(\ell)}{\ell^3}-\frac{3B}{\ell}
 \le\liminf_n\frac{e_{\rm op}(n)}{n^3}
 \le\limsup_n\frac{e_{\rm op}(n)}{n^3}
 \le\frac{e_{\rm op}(\ell)}{\ell^3}+\frac{3B}{\ell}.
\]
The difference of limsup and liminf is at most $6B/\ell$ for every
$\ell$, hence is zero.  This also proves the first finite-size bound.

Closing an open $n$-cube adds $3n^2$ links and their incident
plaquettes.  Adjoin their free electric rotors and use
Lemma~\ref{therm4-cut} to obtain
\[
 |E_{\rm aux}^{\rm per}(n)-e_{\rm op}(n)|\le3Bn^2.
\]
Theorem~\ref{lat-main} identifies the unrestricted periodic minimum
with the physical vacuum minimum: $E_{\rm aux}^{\rm per}=E_0^{\rm per}$.
This proves the remaining assertions.
\end{proof}

The cutting argument concerns the auxiliary tensor product.
It does not factor the Gauss-constrained physical space or presume
that an arbitrary physical boundary flux can be set to zero.

\subsection{A fixed periodic nuclear cell}

Choose a cubic period cell of side $p\ge1$, with $V_c=p^3$ sites,
$K_c\ge1$ labeled positive nuclei, and total charge
$Z_c\le2V_c$.  A rectangular period can be replaced by a cubic
supercell.  Denote the periodically repeated site source by $\nu$.
All periodic source volumes below have $n\in p\mathbb N$ and $n\ge3$.
Choose integers $q_x^{(c)}\in\{0,1,2\}$ with sum $Z_c$ and an
integer flow $F$ on internal cell links such that
\begin{equation}
 \operatorname{div}F=\nu-q^{(c)},\qquad
 \mathcal F_c=\sum_{e\ {\rm in\ cell}}F_e^2<\infty.
 \label{therm4-cell-flow}
\end{equation}
A spanning tree supplies the flow.  For a one-site cell its right-hand
side is zero.  The trial flow has no boundary flux.
Set $\mu=m+3r/a$, $\rho=Z_c/V_c$, and $c_E=g^2/(2a)$.

\begin{proposition}[Two-sided fixed-cell relative bounds]
\label{therm4-cell-bounds}
If $N_c=(n/p)^3$, $K=K_cN_c$, and $Z_{\rm tot}=Z_cN_c$, then
\begin{equation}
 0\le E_\nu^{\rm per}(n)-E_0^{\rm per}(n)-mZ_{\rm tot}
       \le C_{\rm cell}K,
 \label{therm4-cell-relative}
\end{equation}
where one admissible constant is
\begin{equation}
 C_{\rm cell}=
 \frac{3rZ_c+(6+6r+3/g^2)V_c+(g^2/2)\mathcal F_c}{aK_c}.
 \label{therm4-cell-constant}
\end{equation}
Every ground state in this periodic source sector obeys the absolute
electric-energy budget
\begin{equation}
 \langle H_E\rangle\le A_c n^3,\qquad
 A_c=2(h_*-\mu)+\mu\rho
            +c_E\mathcal F_c/V_c+3/(g^2a).
 \label{therm4-electric-budget}
\end{equation}
These constants depend on the fixed cell; they are not estimates
uniform over arbitrary dilute or nonperiodic source arrangements.
\end{proposition}
\begin{proof}
Repeat the cell flow, assign all intercell links electric flux zero,
and occupy two negative-$\beta$ spinors and $q_x^{(c)}$ positive-$\beta$
spinors at each site.  This state satisfies every local Gauss
constraint.  Its hopping expectations vanish.  Its periodic energy is
at most
\begin{equation}
 -2\mu n^3+\mu Z_{\rm tot}
       +c_E N_c\mathcal F_c+3n^3/(g^2a).
 \label{therm4-cell-trial}
\end{equation}
Subtract $E_0^{\rm per}\ge-2h_*n^3$ and $mZ_{\rm tot}$ to obtain the
upper bound.  The lower bound is Corollary~\ref{lat-rest-energy}.
Finally $H_B\ge0$ and $d\Gamma(h_W)\ge-2h_*n^3$ imply
$\langle H_E\rangle\le E_\nu^{\rm per}+2h_*n^3$; insert
\eqref{therm4-cell-trial}.  The budget is absolute and extensive.
It is not an estimate of electric fluctuations relative to the
interacting vacuum proportional to the number of localized defects.
\end{proof}

The same cell trial state on an open union of complete cells gives
\begin{equation}
 E_Y^{\rm op}(\nu)\le -2m|Y|+m\rho|Y|
     +\bigl(c_E\mathcal F_c/V_c+3/(g^2a)\bigr)|Y|.
 \label{therm4-open-source-upper}
\end{equation}
Here each site's fermionic diagonal expectation is
$(q_x^{(c)}-2)(m+r\deg_Y(x)/(2a))\le(q_x^{(c)}-2)m$.
This supplies a uniform extensive bound for all rectangular unions
of whole cells, including boundary remainders.

\begin{theorem}[Open zero-flux source thermodynamic limit]
\label{therm4-open-source-limit}
For every such fixed cell the limit
\begin{equation}
 e_\nu^{\rm op}=\lim_{\substack{n\to\infty\\p\mid n}}
                   \frac{E_{Q_n}^{\rm op}(\nu)}{n^3}
 =\inf_{\substack{\ell\in p\mathbb N\\\ell\ge3}}
       \left\{\frac{E_{Q_\ell}^{\rm op}(\nu)}{\ell^3}
                    +\frac{3B}{\ell}\right\}
 \label{therm4-open-source-limit-eq}
\end{equation}
exists and is finite.  Also
\begin{equation}
 e_{\rm vac}+m\rho
 \le\liminf_{\substack{n\to\infty\\p\mid n}}
                     \frac{E_\nu^{\rm per}(n)}{n^3}
 \le\limsup_{\substack{n\to\infty\\p\mid n}}
                     \frac{E_\nu^{\rm per}(n)}{n^3}
 \le e_\nu^{\rm op}.
 \label{therm4-periodic-bracket}
\end{equation}
In particular the open-boundary relative energy per nucleus has the
limit $(V_c/K_c)(e_\nu^{\rm op}-e_{\rm vac})$.
\end{theorem}
\begin{proof}
Take a product of physical zero-flux block ground states and put the
new crossing rotors in their zero-flux states.  The product satisfies
the source Gauss constraints on the joined graph exactly.  The
bounded coupling estimate gives the one-sided almost-subadditivity
\[
 E_{Q_{k\ell}}^{\rm op}(\nu)
 \le k^3E_{Q_\ell}^{\rm op}(\nu)+3B(k\ell)^3/\ell.
\]
For any $n,\ell\in p\mathbb N$, subdivide as in the vacuum proof.
Every remainder has side lengths divisible by $p$, so it is a union
of admissible complete cells.  Its physical trial energy is
$O(\ell n^2)$ by \eqref{therm4-open-source-upper}.  Thus
\[
 \limsup_{\substack{n\to\infty\\p\mid n}}
 \frac{E_{Q_n}^{\rm op}(\nu)}{n^3}
 \le\frac{E_{Q_\ell}^{\rm op}(\nu)}{\ell^3}+\frac{3B}{\ell}
\]
for every fixed admissible $\ell$.  Let $\gamma$ be the infimum on
the right of \eqref{therm4-open-source-limit-eq}.  By its definition,
$E_{Q_n}^{\rm op}(\nu)/n^3+3B/n\ge\gamma$.  Hence the liminf is
at least $\gamma$ and the preceding limsup is at most $\gamma$.
The extensive lower bound makes this a finite number.
Only upper almost-subadditivity was needed.

Adjoining zero-flux closing rotors to an open physical state gives
$E_\nu^{\rm per}(n)\le E_{Q_n}^{\rm op}(\nu)+3Bn^2$.
The lower bound in \eqref{therm4-periodic-bracket} follows from the
periodic charge-cost theorem and Theorem~\ref{therm4-vacuum-limit}.
\end{proof}

An independently prescribed per-species subtraction
$S_\infty=N_c c_E\mathcal G_\infty(0)\sum_{i\in{\rm cell}}Z_i^2$
simply subtracts its fixed density from these limits and bounds.
The periodic prescription $S_n$ has the same limiting density.
Indeed $\mathcal G_n(0)\to\mathcal G_\infty(0)$: away from the origin
the Fourier sum is a Riemann sum, while on $0<|k|\le\delta n$
the shell argument in Lemma~\ref{lat-green-bound} bounds its
contribution by $C\delta$ uniformly in $n$; the corresponding integral
has the same integrable $|t|^{-2}$ control.  First let $n\to\infty$
off this ball and then let $\delta\downarrow0$.

\subsection{The exact boundary-flux reduction}

For a partition into blocks, orient every cut link once.  For
$f\in\mathbb Z^{\mathcal C}$ define the outward flux function at a
boundary site of $B$ by
\[
 b_B(f)_x=\sum_{\substack{e\in\mathcal C\\e\ {\rm leaves}\ B\ {\rm at}\ x}}
 f_e-\sum_{\substack{e\in\mathcal C\\e\ {\rm enters}\ B\ {\rm at}\ x}}f_e.
\]
Inside this block Gauss's law becomes
$\operatorname{div}_BE+q=\nu|_B-b_B(f)$.
Set a sector energy to $+\infty$ if its sector is empty.

\begin{proposition}[Cutting identity for constrained minima]
\label{therm4-flux-identity}
With the preceding conventions,
\begin{equation}
 \left|E_\nu-
 \inf_{f\in\mathbb Z^{\mathcal C}}
 \left\{\sum_B E_B^{\rm op}(\nu|_B-b_B(f))
                         +c_E\sum_{e\in\mathcal C}f_e^2\right\}
 \right|\le Bn_c.
 \label{therm4-flux-identity-eq}
\end{equation}
Only fluxes satisfying
$|\sum_{x\in B}(\nu_x-b_B(f)_x)|\le2|B|$ in each individual block contribute
finite block energies.
\end{proposition}
\begin{proof}
The cut Hamiltonian of Lemma~\ref{therm4-cut} still commutes with
every original Gauss generator.  Its cut rotors have a simultaneous
integer electric basis.  At fixed $f$, its physical subspace is the
tensor product of the block sectors with sources $\nu|_B-b_B(f)$,
and its energy is the sum displayed in \eqref{therm4-flux-identity-eq}.
The direct sum over $f$ proves the formula for the cut minimum.
The bounded operator $V_{\mathcal C}$ also preserves the original
Gauss sector, so the variational comparison proves the inequality.
\end{proof}

Thus cutting a source-sector state produces different block sources,
not necessarily the zero-flux sectors used in
Theorem~\ref{therm4-open-source-limit}.  Vacuum domination bounds such a
block by its vacuum energy and net-charge cost.  Combined with a cell
upper bound, this loses an error proportional to the block volume;
it does not supply the boundary estimate needed below.

\begin{proposition}[Admissibility obstruction to an unrestricted
boundary-flux lemma]
\label{therm4-admissibility-obstruction}
An inequality comparing every $E_B^{\rm op}(\nu-b_B(f))$ with
$E_B^{\rm op}(\nu)$ at a finite boundary cost is false unless
admissibility of the zero-flux comparison sector is imposed.
This failure occurs for a periodic source of admissible mean density
and individual nuclear charge four.
\end{proposition}
\begin{proof}
Take $\nu_x=4$ on even $x_1$ layers and zero on odd layers.
Its mean is two per site, so a compatible even torus is admissible.
Take a cube of odd side $\ell$ starting on an even layer.  Its source
charge is $2\ell^3+2\ell^2>2|B|$, so $E_B^{\rm op}(\nu)=+\infty$.
On the other hand choose $q_x=2$ everywhere, and internal electric
fields $E_{x,1}=(-1)^{x_1}$, with the other components zero.
The required exterior outward flux is $+1$ on both $x_1$ faces
and zero on the other faces.  Its sum is $2\ell^2$ and
$\nu-b_B(f)$ has total charge $2|B|$.  The displayed occupation
and internal electric eigenvector satisfy Gauss's law, and have
finite energy.  No finite multiple of boundary area and
$\sum f_e^2$ repairs the claimed comparison with $+\infty$.
\end{proof}

Commensurate boxes containing complete period cells avoid this
particular obstruction.  The example does not disprove a suitably
restricted boundary-screening inequality for them.

\subsection{What averaging proves, and what it still requires}

The following is a precise sufficient condition, not an established
property of the Wilson Hamiltonian.
For every commensurate cube $B$ of side $\ell$, suppose constants
$0\le C_\partial,C_f<\infty$, independent of $\ell$ and of integer boundary
flux, obey
\begin{equation}
 E_B^{\rm op}(\nu-b_B(f))
 \ge E_B^{\rm op}(\nu)
       -C_\partial|\partial_{\rm links}B|
       -C_f\sum_{e\in\partial_{\rm links}B}f_e^2.
 \tag{BF}\label{therm4-boundary-hypothesis}
\end{equation}
Here $\partial_{\rm links}B$ counts external link incidences, and
the left sector is assumed admissible.  At a shared cut link each
of its two incidences occurs in this convention.

\begin{proposition}[Conditional periodic-source limit; cut averaging]
\label{therm4-conditional-periodic}
If \eqref{therm4-boundary-hypothesis} holds for the fixed cell, then
\[
 \lim_{\substack{n\to\infty\\p\mid n}}
       E_\nu^{\rm per}(n)/n^3=e_\nu^{\rm op}.
\]
It suffices that $C_f$ be finite; the restriction $C_f<c_E$ is
unnecessary when averaged cut energies are used.
\end{proposition}
\begin{proof}
First let $n$ be divisible by a fixed commensurate $\ell$.
Translate the cut grid by multiples of $p$.  A given link is cut with
probability at most $p/\ell$.  For a periodic source ground state,
\eqref{therm4-electric-budget} therefore gives a grid with
\begin{equation}
 \left\langle\sum_{e\in\mathcal C}E_e^2\right\rangle
 \le\frac{pA_c}{c_E\ell}n^3.
 \label{therm4-cut-average}
\end{equation}
All resulting cubes have the same period-cell phase and
$n_c=3n^3/\ell$.  Apply (BF) in the cut-flux decomposition of the
state.  The sum of boundary incidences is $2n_c$, and a squared
cut flux is also counted twice.  Restoring the bounded cut terms
gives
\begin{equation}
 \frac{E_\nu^{\rm per}(n)}{n^3}
 \ge\frac{E_{Q_\ell}^{\rm op}(\nu)}{\ell^3}
 -\frac{3B+6C_\partial}{\ell}
 -\frac{(2C_f-c_E)_+pA_c}{c_E\ell}.
 \label{therm4-conditional-bound}
\end{equation}
If one instead absorbs the boundary-flux term without averaging,
the incidence convention requires $2C_f\le c_E$, not merely
$C_f<c_E$.

For arbitrary $n\in p\mathbb N$, use full $\ell$-cubes and
rectangular remainders, all with side lengths divisible by $p$.
Average the origin over all translations by complete cells.
In each direction there are $\lceil n/\ell\rceil$ cut planes, so
the probability that a link is cut is at most
$p/\ell+p/n$.  The leftover site count is $O(\ell n^2)$.
Use (BF) only in full cubes and the lower bound
$H_Y^{\rm op}\ge-2h_*|Y|$ in the leftover blocks.  This adds a
vanishing error after division by $n^3$ and passage $n\to\infty$
at fixed $\ell$.  Thus \eqref{therm4-conditional-bound} holds for
the liminf along all compatible sizes.
Let $\ell\to\infty$ through multiples of $p$ and use
Theorem~\ref{therm4-open-source-limit}; its periodic upper bound
supplies the matching limsup.
\end{proof}

The elementary averaging estimate \eqref{therm4-cut-average} is
unconditional.  It controls how much electric energy is exposed by
a cut, but it does not change a block's Gauss sector.  Inequality
(BF), or another argument with the same boundary-scale consequence,
remains unproved here.

The Coulomb thermodynamic framework of
Hainzl--Lewin--Solovej~\cite{HLSThermoI,HLSThermoII} has different
localization hypotheses. The Graf--Schenker inequality stated in
\cite[Theorem~2]{HLSThermoII} averages localized continuum point-charge
$1/|x-y|$ energies over translated and rotated simplices, with an
error of order $\ell^{-1}\sum_i z_i^2$.  Their periodic-crystal
application has nonrelativistic quantum electrons and static
classical nuclei.  It does not include compact electric rotors or
the integer-flux physical-sector decomposition above.  Eliminating
the longitudinal energy in this model leaves a charge-dependent
allowed transverse-flux space, as already noted after
\eqref{lat-longitudinal-split}.  Consequently the continuum
electrostatic localization inequality cannot simply replace (BF).
An appropriate lattice screening or gluing theorem would be new
input to the present proof.

We have therefore proved a physical vacuum limit, two-sided fixed-cell
bounds, and an open zero-flux source limit.  Equality with a periodic
source limit is conditional on the explicitly stated boundary
comparison.  None of these fixed-$a$ results supplies a continuum
limit, binding, a neutral excitation gap, or crystalline order.
